\scanpage{57}
d'un graphe cubique $\Delta$ est isomorphe au groupe $\unsure{\mathfrak S_3|\mathscr S_3}$, avec $-$ correspondant à $\sigma_{12}$, $\sim$ à $\sigma_{23}$ ....

\uline{NB} En \unsure{termes|terme} d'un couple $(X,\{\unsure{\tau_i|\tau_i'}\}_{i\in I})$ d'un bitriangle $X\subset\Delta$ et d'une coorientation de celui-ci, \unsure{les règles|donc presque} précédentes associent à un tel couple, d'une part le
\deleted{couple $(X',(\tau_i')_{i\in I})$ \inserted{\unsure{même} $I$} formé du bitriangle $X'$ associé à $X$, tel que les $\tau_i'$ sont les orbites}
$X'=\bigcup_{i\in I}\overline{\tau_i'}$ \inserted{triade opposée} muni des $(\overline{\tau_i'})_{i\in I}$, d'autre part $X$ muni de l'autre coorientation $(\tau_i')_{i\in I'}$ .....

\deleted{\textcircled{15} \unsure{Pour les triades}}

\textcircled{15} Les triangles de $X_\alpha$ correspondent aux parties $\deleted{\unsure{h}} \inserted{t}\subset\prod_{\beta\in\Lambda-\{\alpha\}} I_\beta$ \deleted{qui sont des} telles que les applications $t\longrightarrow I_\beta$ $(\beta\in\unsure{\Delta|\Lambda}-\{\alpha\})$ \inserted{i.e. (si $\Lambda=\{\alpha,\beta,\gamma\}$) aux bijections $I_\beta\overset{\sim}{\longrightarrow}I_\gamma$} sont bijectives, en associant à une telle \deleted{partie} \deleted{\unsure{l'ens des él.}} \inserted{bijection $\varphi\colon\deleted{\unsure{I}}$} l'ens des él. $\{\{x_\beta,\varphi(x_\beta)\}\}_{x_\beta\in I_\beta}$ de $X_\alpha$. Ceci posé, pour un tel triangle $\underline{t}=\underline{t}_\varphi$, la symétrie correspondante de $\Delta$ est

% archive sheet number: 56

